\section{Related Work}
\label{sec:related_work}

%Over the years, deep learning has revolutionized the medical image segmentation domain by providing SOTA performance on a wide range of tasks. This section reviews the advancements in CNN-based, Vision Transformer-based, and Lightweight models in medical image segmentation, that are relevant to our proposed MK-UNet.

\subsection{Convolutional Neural Networks (CNNs)}
The advent of CNNs marks a significant shift in medical image segmentation \cite{ronneberger2015u, oktay2018attention, zhou2018unet++, fan2020pranet, kim2021uacanet, chen2017deeplab}. Pioneering work such as Fully Convolutional Networks (FCNs) \cite{long2015fully} laid the foundation for end-to-end segmentation models. FCNs replace fully connected layers with convolutional layers, thus enabling pixel-wise predictions and efficient learning of spatial hierarchies in images. U-Net \cite{ronneberger2015u} became a cornerstone in medical image segmentation due to its encoder-decoder architecture with skip connections. This design effectively combines high-resolution features from the encoder with context information from the decoder, hence leading to precise segmentation even with limited training data. The sophisticated design of U-shaped architecture for pixel-level segmentation tasks motivates us to choose the U-shaped design in our proposed network.

U-Net's success has inspired numerous variants and improvements. Inspired by residual learning in ResNet \cite{he2016deep}, ResUNet \cite{zhang2018road} uses residual blocks to facilitate gradient flow and improve convergence, addressing the vanishing gradient problem in deep networks. Zhou et al.\cite{zhou2018unet++} introduce UNet++, which uses dense nested skip connections to further enhance the feature propagation and improve the segmentation accuracy. AttnUNet \cite{oktay2018attention} incorporates attention mechanisms to focus on relevant regions in the feature maps, improving segmentation performance by suppressing irrelevant background noise. Fan et al. \cite{fan2020pranet} introduce PraNet for precise polyp segmentation that employs parallel reverse attention and edge-guidance to refine segmentation boundaries. UACANet \cite{kim2021uacanet} uses uncertainty-aware mechanisms to improve the reliability and robustness of segmentation outcomes. DeepLabv3+ \cite{chen2017deeplab} integrates atrous convolutions and spatial pyramid pooling to capture multi-scale context information. ACC-UNet \cite{ibtehaz2023acc} employs adaptive context capture mechanisms to dynamically adjust the receptive fields based on the input image.

\subsection{Vision Transformers}
Vision Transformers (ViTs) \cite{dosovitskiy2020image, liu2021swin} have emerged as a powerful alternative to CNNs, offering a new paradigm for medical image analysis tasks by leveraging the self-attention mechanism \cite{chen2021transunet, cao2021swin, rahman2024emcad, Rahman_2023_WACV, rahman2023multi, rahman2024g, valanarasu2021medical}. By combining the strengths of CNNs for local feature extraction and Transformers for capturing long-range dependencies, TransUNet \cite{chen2021transunet} achieves superior performance in medical image segmentation. SwinUNet \cite{cao2021swin} is introduced based on the Swin Transformer \cite{liu2021swin} architecture, which utilizes shifted windows to achieve hierarchical feature representation, enabling efficient computation. MedT \cite{valanarasu2021medical}, a lightweight Transformer model specifically designed for medical image segmentation, which employs gated axial attention mechanisms to focus on relevant regions and reduce computational complexity. Rahman et al. introduce CASCADE \cite{Rahman_2023_WACV}, a cascaded-attention decoding network using standard convolutions. Recently, EMCAD \cite{rahman2024emcad} introduces a depthwise convolutions-based multi-scale decoder. Although, CASCADE and EMCAD perform well in medical image segmentation, their segmentation accuracy and computational complexity solely depend on the strength and complexity of the exiting pretrained transformer encoder they use, thus making them less suitable for resource-constrained settings. In contrast, we propose to design an extremely efficient (ultra-lightweight) end-to-end (both encoder and decoder) architecture using the multi-kernel trick (where, $k_1=k_2$ or $k_1 \neq k_2$ for $k_1, k_2 \in Kernels$) with depthwise convolutions.

\subsection{Lightweight CNNs}
Recent efforts have focused on making CNNs more efficient for real-time and resource-constrained environments. MobileNets \cite{howard2017mobilenets} and EfficientNets \cite{tan2019efficientnet} introduce depthwise separable convolutions and compound scaling, respectively, to create lightweight models with competitive performance. Additionally, several novel lightweight architectures have been developed to further enhance the efficiency of medical image segmentation \cite{valanarasu2022unext, tang2023cmunext, ruan2023ege, liu2024rolling}. UNeXt \cite{valanarasu2022unext} uses hybrid convolutional and transformer blocks to capture local and global features efficiently, improving segmentation accuracy while maintaining computational efficiency. CMUNeXt \cite{tang2023cmunext} combines convolutional and multi-scale features to enhance segmentation performance. EGE-UNet \cite{ruan2023ege} integrates edge-guided mechanisms to refine segmentation boundaries. Rolling-UNet \cite{liu2024rolling} incorporates rolling convolutional blocks to enhance the model's ability to capture long-range dependencies.